% GENERATED FILE. Edit canonical lessons, then run npm run generate:books.
% canonical-source-hash: fnv1a64:b888b5a4d98b229c
% canonical-lessons: RU-C378-dogovarivatsya, RU-C378-dogovoryonnost, RU-C378-variant, RU-C378-alternativa, RU-C378-okonchatelnyy

\chapter{Making arrangements}
\label{ch:ru-making-arrangements}
\hlchaptermodality{378}

\begin{chapteropening}
\textbf{By the end of this chapter:} \emph{I can make an arrangement with someone, compare options and alternatives, and say that a decision is final.}

Complete the last lesson of chapter 378: I can make an arrangement with someone, compare options and alternatives, and say that a decision is final.
\end{chapteropening}

\section[\texorpdfstring{\ru{договариваться} (dogovárivatsya)}{договариваться (dogovárivatsya)}]{\ru{договариваться} (dogovárivatsya) — to arrange, to come to an agreement}

\label{lesson:RU-C378-dogovarivatsya}



\begin{quote}
Before the new one: say the Russian for to set, then the Russian for to be going to.
\end{quote}

\subsection*{You'll want to know: \ru{договариваться}}
\textbf{\ru{договариваться}} (\emph{dogovárivatsya}) — \textquotedblleft{}to arrange, to come to an agreement\textquotedblright{}.

\begin{quote}
\ru{Мы} \ru{договорились} \ru{встретиться} \ru{у} \ru{метро}.
\end{quote}

\begin{quote}
\emph{My dogovorílis vstrétitsya u metró.}
\end{quote}

\begin{quote}
We agreed to meet by the metro.
\end{quote}

Imperfective; perfective \emph{\ru{договориться}}. \emph{\ru{Договариваться} \ru{с} \ru{кем} \ru{о} \ru{чём}}, or + infinitive.

\begin{cousinweb}
\textbf{\ru{договариваться}} is built from \emph{\ru{до}-} \textquotedblleft{}through to the end\textquotedblright{} and \emph{\ru{говорить}} \textquotedblleft{}to speak\textquotedblright{}, on the old Slavic noun \emph{*govorŭ} \textquotedblleft{}talk\textquotedblright{} (EtymDB): to talk a matter through.
\end{cousinweb}

\begin{grammarlens}[title={What you've built}]
One more word for telling, explaining and describing at B1.
\end{grammarlens}

\subsection*{Your turn}
\begin{itemize}
\raggedright
  \item \emph{Say it:} \emph{dogovárivatsya}
  \item \emph{Say it:} \emph{dogovárivatsya}, once more
  \item \emph{Say it:} a sentence of your own with \emph{dogovárivatsya}
  \item \emph{Recall:} say the Russian for to set, then the Russian for to be going to, then say \emph{dogovárivatsya} again
\end{itemize}

\begin{tcolorbox}[breakable,colback=teal!4,colframe=teal!35!black,title={Before you move on}]
What does \ru{договариваться} mean? (\textquotedblleft{}To arrange, to come to an agreement\textquotedblright{}.) Say it once more.
\end{tcolorbox}
\section[\texorpdfstring{\ru{договорённость} (dogovoryónnost)}{договорённость (dogovoryónnost)}]{\ru{договорённость} (dogovoryónnost) — an arrangement, an understanding}

\label{lesson:RU-C378-dogovoryonnost}



\begin{quote}
Before the new one: say the Russian for to be going to, then the Russian for to arrange.
\end{quote}

\subsection*{You'll want to know: \ru{договорённость}}
\textbf{\ru{договорённость}} (\emph{dogovoryónnost}) — \textquotedblleft{}an arrangement, an understanding\textquotedblright{}.

\begin{quote}
\ru{У} \ru{нас} \ru{есть} \ru{договорённость} \ru{с} \ru{директором}.
\end{quote}

\begin{quote}
\emph{U nas yest dogovoryónnost s diréktorom.}
\end{quote}

\begin{quote}
We have an arrangement with the director.
\end{quote}

Feminine. An arrangement in words; a written agreement is \emph{\ru{договор}}.

\begin{cousinweb}
\textbf{\ru{договорённость}} is built on the participle of \emph{\ru{договориться}} \textquotedblleft{}to agree\textquotedblright{}, from \emph{\ru{говорить}} \textquotedblleft{}to speak\textquotedblright{}, with the abstract suffix \emph{-\ru{ость}}.
\end{cousinweb}

\begin{grammarlens}[title={What you've built}]
One more word for telling, explaining and describing at B1.
\end{grammarlens}

\subsection*{Your turn}
\begin{itemize}
\raggedright
  \item \emph{Say it:} \emph{dogovoryónnost}
  \item \emph{Say it:} \emph{dogovoryónnost}, once more
  \item \emph{Say it:} a sentence of your own with \emph{dogovoryónnost}
  \item \emph{Recall:} say the Russian for to be going to, then the Russian for to arrange, then say \emph{dogovoryónnost} again
\end{itemize}

\begin{tcolorbox}[breakable,colback=teal!4,colframe=teal!35!black,title={Before you move on}]
What does \ru{договорённость} mean? (\textquotedblleft{}An arrangement, an understanding\textquotedblright{}.) Say it once more.
\end{tcolorbox}
\section[\texorpdfstring{\ru{вариант} (variánt)}{вариант (variánt)}]{\ru{вариант} (variánt) — an option, a version}

\label{lesson:RU-C378-variant}



\begin{quote}
Before the new one: say the Russian for to arrange, then the Russian for an arrangement.
\end{quote}

\subsection*{You'll want to know: \ru{вариант}}
\textbf{\ru{вариант}} (\emph{variánt}) — \textquotedblleft{}an option, a version\textquotedblright{}.

\begin{quote}
\ru{Есть} \ru{два} \ru{варианта}: \ru{поезд} \ru{или} \ru{самолёт}.
\end{quote}

\begin{quote}
\emph{Yest dva variánta: póyezd íli samolyót.}
\end{quote}

\begin{quote}
There are two options: the train or the plane.
\end{quote}

Masculine. \emph{\ru{Лучший} \ru{вариант}} — the best option. It is also a version of a text.

\begin{cousinweb}
\textbf{\ru{вариант}} was borrowed from French \emph{variante} (EtymDB), which goes back to Latin \emph{variāns}, the participle of \emph{variāre} \textquotedblleft{}to change, to vary\textquotedblright{}, built on \emph{varius} \textquotedblleft{}various\textquotedblright{} (Lewis \& Short).
\end{cousinweb}

\begin{grammarlens}[title={What you've built}]
One more word for telling, explaining and describing at B1.
\end{grammarlens}

\subsection*{Your turn}
\begin{itemize}
\raggedright
  \item \emph{Say it:} \emph{variánt}
  \item \emph{Say it:} \emph{variánt}, once more
  \item \emph{Say it:} a sentence of your own with \emph{variánt}
  \item \emph{Recall:} say the Russian for to arrange, then the Russian for an arrangement, then say \emph{variánt} again
\end{itemize}

\begin{tcolorbox}[breakable,colback=teal!4,colframe=teal!35!black,title={Before you move on}]
What does \ru{вариант} mean? (\textquotedblleft{}An option, a version\textquotedblright{}.) Say it once more.
\end{tcolorbox}
\section[\texorpdfstring{\ru{альтернатива} (alternatíva)}{альтернатива (alternatíva)}]{\ru{альтернатива} (alternatíva) — an alternative}

\label{lesson:RU-C378-alternativa}



\begin{quote}
Before the new one: say the Russian for an arrangement, then the Russian for an option.
\end{quote}

\subsection*{You'll want to know: \ru{альтернатива}}
\textbf{\ru{альтернатива}} (\emph{alternatíva}) — \textquotedblleft{}an alternative\textquotedblright{}.

\begin{quote}
\ru{У} \ru{нас} \ru{нет} \ru{альтернативы}.
\end{quote}

\begin{quote}
\emph{U nas net alternatívy.}
\end{quote}

\begin{quote}
We have no alternative.
\end{quote}

Feminine. \emph{\ru{Альтернатива}} + dative, or \emph{\ru{для}}: \emph{\ru{альтернатива} \ru{машине}}.

\begin{cousinweb}
\textbf{\ru{альтернатива}} is usually traced to French \emph{alternative}, built on Latin \emph{alternus} \textquotedblleft{}by turns\textquotedblright{}, from \emph{alter} \textquotedblleft{}the other of two\textquotedblright{} (Lewis \& Short).
\end{cousinweb}

\begin{grammarlens}[title={What you've built}]
One more word for telling, explaining and describing at B1.
\end{grammarlens}

\subsection*{Your turn}
\begin{itemize}
\raggedright
  \item \emph{Say it:} \emph{alternatíva}
  \item \emph{Say it:} \emph{alternatíva}, once more
  \item \emph{Say it:} a sentence of your own with \emph{alternatíva}
  \item \emph{Recall:} say the Russian for an arrangement, then the Russian for an option, then say \emph{alternatíva} again
\end{itemize}

\begin{tcolorbox}[breakable,colback=teal!4,colframe=teal!35!black,title={Before you move on}]
What does \ru{альтернатива} mean? (\textquotedblleft{}An alternative\textquotedblright{}.) Say it once more.
\end{tcolorbox}
\section[\texorpdfstring{\ru{окончательный} (okonchátelnyy)}{окончательный (okonchátelnyy)}]{\ru{окончательный} (okonchátelnyy) — final, definitive}

\label{lesson:RU-C378-okonchatelnyy}



\begin{quote}
Before the new one: say the Russian for an option, then the Russian for an alternative.
\end{quote}

\subsection*{You'll want to know: \ru{окончательный}}
\textbf{\ru{окончательный}} (\emph{okonchátelnyy}) — \textquotedblleft{}final, definitive\textquotedblright{}.

\begin{quote}
\ru{Это} \ru{моё} \ru{окончательное} \ru{решение}.
\end{quote}

\begin{quote}
\emph{Éto moyó okonchátelnoye reshéniye.}
\end{quote}

\begin{quote}
This is my final decision.
\end{quote}

\emph{\ru{Окончательно}} — finally, for good.

\begin{cousinweb}
\textbf{\ru{окончательный}} is built on \emph{\ru{окончить}} \textquotedblleft{}to finish\textquotedblright{}, from \emph{\ru{конец}} \textquotedblleft{}an end\textquotedblright{}, the root that also gives \emph{\ru{конечно}} \textquotedblleft{}of course\textquotedblright{}.
\end{cousinweb}

\begin{grammarlens}[title={What you've built}]
One more word for telling, explaining and describing at B1.
\end{grammarlens}

\subsection*{Your turn}
\begin{itemize}
\raggedright
  \item \emph{Say it:} \emph{okonchátelnyy}
  \item \emph{Say it:} \emph{okonchátelnyy}, once more
  \item \emph{Say it:} a sentence of your own with \emph{okonchátelnyy}
  \item \emph{Recall:} say the Russian for an option, then the Russian for an alternative, then say \emph{okonchátelnyy} again
\end{itemize}

\begin{tcolorbox}[breakable,colback=teal!4,colframe=teal!35!black,title={Before you move on}]
What does \ru{окончательный} mean? (\textquotedblleft{}Final, definitive\textquotedblright{}.) Say it once more.
\end{tcolorbox}
